\documentclass[10pt,a4paper]{amsart}
\usepackage[margin=19mm]{geometry}
\usepackage{amsmath,amssymb,mathtools}
\usepackage[scr=rsfs,cal=euler]{mathalfa}
\usepackage{xfrac}
\usepackage[unicode,hidelinks,bookmarks=false]{hyperref}
\newcommand{\ie}{i.e.\ }
\newcommand{\cf}{cf.\ }
\newcommand{\resp}{respectively\ }
\newcommand{\Subsection}{\subsection}
\providecommand{\nobreakdash}{\nobreak\hbox{-}}
\setlength{\emergencystretch}{2em}
\multlinegap0pt
\usepackage{amssymb}
\usepackage{enumerate}
\usepackage{mathtools}
\newtheorem{lemma}{Lemma}
\newtheorem{proposition}{Proposition}
\newcommand\NN{\mathbb{N}}
\newcommand\ZZ{\mathbb{Z}}
\renewcommand{\c}{\mathbf{c}}
\renewcommand{\d}{\mathrm{d}}
\newcommand{\fp}{\mathfrak{p}}
\renewcommand{\geq}{\geqslant}
\renewcommand{\leq}{\leqslant}
\DeclareMathOperator{\n}{N}
\DeclareMathOperator{\nf}{\mathbf{N}}
\newcommand{\ve}{\varepsilon}
\newcommand{\x}{\mathbf{x}}
\title{Bateman--Horn, polynomial Chowla and the Hasse principle with probability 1}
\author{Tim Browning, Efthymios Sofos, Joni Ter\"av\"ainen}
\date{Source: arXiv:2212.10373v2 (CC BY 4.0)}
\begin{document}
\maketitle

\noindent\emph{Source and licence.} Bateman--Horn, polynomial Chowla and the Hasse principle with probability 1. Version arXiv:2212.10373v2 (CC BY 4.0), distributed by arXiv under the Creative Commons Attribution 4.0 International licence, http://creativecommons.org/licenses/by/4.0/, as recorded in the arXiv licence metadata for this exact version and shown on the abstract page https://arxiv.org/abs/2212.10373v2. Full licence notice: \texttt{licenses/arxiv-2212.10373-CC-BY-4.0.txt}.\par\smallskip

\noindent\textit{Editorial scope.} Counted unit: the article's lemma lem:tilde-N (source0.tex lines 4088--4110). The article's complete original proof of it is delivered with it (source0.tex lines 4112--4402, one \texttt{proof} environment of 291 lines). No mathematics is written, repaired or re-derived: the statement and the proof are the article's own lines, in the article's own notation, and no retained line is substituted or re-flowed.  Retained uncounted as the source-local context the proof needs: lines 1884--1886; lines 1891--1901; lines 1918--1925; lines 2377--2379; lines 2381--2384; lines 2578--2591; lines 2606--2608; lines 2632--2634; lines 2636--2638; lines 2655--2675; lines 2683--2688; lines 2790--2792; lines 2795--2801; lines 2943--2949; lines 2959--2962; lines 3695--3698; lines 3973--3976; lines 4004--4008.  Nothing was omitted from any retained span, so the package carries no omission record.  No retained line contains a citation, so the wrapper carries no bibliography and no bibliography entry.  No retained line contains a figure environment, an image-inclusion command or drawing code, so no image is delivered and the package declares no asset dependency.  Editorial formatting only, in addition: an AMS \texttt{amsart} preamble that restates the article's own declarations and macros used by the retained lines, namely the environments \texttt{lemma}, \texttt{proposition} on separate counters, together with the standard packages those lines need.  The retained environments print in this document as Lemma 1, Proposition 1 in order of appearance, whereas the article prints the same items with the numbering of its own full document.  The complete native source of the article as distributed in the arXiv e-print for this exact version is bundled unchanged as \texttt{sources/135134/source0.tex}.
\begin{equation}\label{eq:lambda}
\lambda_f(q)=\#\left\{u\in \ZZ/q\ZZ: f(u)\equiv 0\bmod{q}\right\},
\end{equation}

\begin{lemma}\label{lem:stewart}
Let $p$ be a prime and let $k\in \NN$. Let $f\in \mathbb{Z}[t]$ be a polynomial of degree $d\geq 1$. Assume that $f$ is separable  and let $\sigma$ be the $p$-adic valuation of the  content of $f$. 
If $\sigma\geq k$ then $\lambda_f(p^k)=p^k$. If $\sigma<k$ then
$$
\lambda_f(p^k)\leq d \min\left\{p^{k(1-\frac{1}{d})+\frac{\sigma}{d}},  p^{k-1}\right\}.
$$
Additionally, if $p$ does not divide the discriminant of $f$, then
\begin{align*}
\lambda_f(p^k)\leq d.     
\end{align*}
\end{lemma}

\begin{lemma}\label{lem:henriot}
Let $\delta>0$ and $A,C>0$. Let $f\in \ZZ[t]$ be a 
separable 
polynomial of degree $d\geq 1$. Let $\|f\|$ be the maximum modulus of the coefficients of $f$ and assume that the content of $f$ is at most $(\log x)^A$. Then there exist a positive constant $K$, depending only on $\delta$, $d$ and $C$, such that for $x\geq \|f\|^{\delta}$ we have
$$
\sum_{\substack{n\leq x}} \tau(|f(n)|)^C \ll x (\log x)^{K}.
$$
\end{lemma}

\begin{equation}\label{eq:k-x}
k(x)=1000dA \lceil \log\log x\rceil
\end{equation}

\begin{equation}\label{eq:def_W}
W=\prod_{p\leq w} p^{k(x)}\quad \text{ for }\quad 
w=\exp(\sqrt{\log x}).
\end{equation}

\begin{proposition}\label{prop:1moment} Recall the notation~\eqref{eq:k-x} and~\eqref{eq:def_W}.
Let $X$ satisfy  $B^e\geq X\geq x$ and $W\geq X^2$. Let $I\subset [1,X]$ such that $|I|>X^{9/10}$. 
Let $u,q\in \NN$ such that $1\leq u\leq q\leq X^{1/5}$. Assume that $p^{k(x)}\nmid q$ for any $p\leq w$, and that
 \begin{equation}\label{eq:assume-gcd}
\gcd(u,q)\leq \exp \left(\frac{2\sqrt{\log X}}{\log\log X}\right).
\end{equation}
Then 
\begin{align*}
S(I;q,u)=~& \frac{\gamma(\gcd(q,W),u)}{q}\int_{I} \omega(t;B)\d t  +
 O\left(\frac{|I|}{q}
 \exp\left(-\frac{\log X}{23\sqrt{\log x}} \right)
\right).
\end{align*}
\end{proposition}

\begin{equation}\label{eq:W0}
W_0\ll (\log x)^{C'}\leq (\log X)^{C'},
\end{equation}

\begin{equation}\label{eq:alpha-p}
\alpha_p=\left(1-\frac{1}{p}\right)^{-1}\prod_{\fp\mid p} \left(1-\frac{1}{\n \fp}\right)
\end{equation}

\begin{equation}\label{eq:beta-p}
\beta_{p^k}=p^k\sum_{j\geq k} \frac{r_K(p^j)}{p^j}\prod_{\fp\mid p} \left(1-\frac{1}{\n \fp}\right).
\end{equation}

\begin{equation}\label{eq:calc}
\begin{split}
\gamma(W_1,n)
&=\prod_{\substack{p\mid P(w)}} \gamma(p^{k(x)},n) \\
&=\prod_{\substack{p\mid P(w)\\
p^{k(x)}\nmid n}}
r_K(p^{v_p(n)}) \alpha_p
\prod_{\substack{p\mid P(w)\\
p^{k(x)}\mid n
}}
\beta_{p^{k(x)}}\\
&=c(w)
\prod_{\substack{p\mid P(w)\\
p^{k(x)}\nmid n}}
r_K(p^{v_p(n)}) 
\prod_{\substack{p\mid P(w)\\
p^{k(x)}\mid n
}}
\alpha_p^{-1}\beta_{p^{k(x)}},
\end{split}
\end{equation}

\begin{equation}\label{eq:r-tilde}
\tilde r_K (p^j) =\begin{cases}
r_K(p^j) &\text{ if $j<k(x)$,}\\
\alpha_p^{-1}\beta_{p^{k(x)}} &\text{ if $j\geq k(x)$.}
\end{cases}
\end{equation}

\begin{equation}\label{eq:definition-sig0}
\Sigma_B(z)=\sum_{\substack{p\mid k \Rightarrow p\leq w\\ k>z}} \frac{\tau(k)^{B}}{k},
\end{equation}

\begin{lemma}\label{lem:gouda}
Let $B\geq 1$, $w\geq 2$ and $z\geq w^{C_B}$, where $C_B$ is a large enough constant. Then 
$$
\Sigma_{B}(z)\ll_B
\exp\left(-\frac{\log z}{4\log w} \right).
$$
\end{lemma}

\begin{equation}\label{eq:nature-g}
|g(p^j)|
\begin{cases}
=0 &\text{ if $j>k(x)$,}\\
\leq  \tau(p^j)^{e+1} &\text{ if $j\leq k(x)$.}
\end{cases}
\end{equation}

\begin{equation}\label{eq:size-g}
|g(k)|=\prod_{p^j\| k} |g(p^j)| \leq 
 \tau(k)^{e+1} \quad \text{ if }\, k\in \mathcal{C}_w.
\end{equation}

\begin{equation}\label{eq:Sd+}
S_d(\tilde H, H)=
\left\{\c\in S_d(H): \frac{1}{2}\tilde H\leq c_0\leq \tilde H \right\}
\end{equation}

\begin{equation}\label{eq:WWW'}
W_0=\prod_{\substack{\text{$p\leq M_{d,K}$ or $p\mid h_\c$}}} p^{k(x)} \quad \text{ and } \quad 
W_1=\prod_{\substack{M_{d,K}<p\leq w \\ p\nmid h_\c}}p^{k(x)} .
\end{equation}

\begin{equation}\label{eq:sig-c}
\sigma_\c(W_0)=W_0^{-e}\#\left\{
(\x,t)\in (\ZZ/W_0\ZZ)^{e+1}: \nf_K(\x)\equiv f_\c(t)\bmod{W_0}
\right\}.
\end{equation}

\begin{lemma}\label{lem:tilde-N}
Let $A\geq 1$.
Let $\c\in S_d(\tilde H, H)$, in the notation of~\eqref{eq:Sd+}. 
Recall the definition ~\eqref{eq:r-tilde} of 
$\tilde r_K$.
Suppose that $f_\c$ is separable, has no integer zero, and has content $h_\c\leq (\log x)^A$.
Then
$$
\tilde N_\c(I) = 
\sigma_\c(W_0)|I| 
\sum_{\substack{k\in \NN\\ p\mid k \Rightarrow p\mid P_\c(w) }}
\frac{c_\c(w) g(k) 	
\lambda_{f_\c}(k)}{k}
+O_A\left(
x\exp\left(-(\log x)^{1/8}\right)\right),
$$
where $g=\tilde r_K*\mu$, $\lambda_{f_\c}$ is given by~\eqref{eq:lambda}, 
 $W_0$ is given by ~\eqref{eq:WWW'} and 
$$
c_\c(w)=\prod_{p\mid P_\c(w)} \alpha_p,
$$
with  $\alpha_p$  given by~\eqref{eq:alpha-p}.
\end{lemma}

\begin{proof}
We shall use 
similar techniques to those that we developed in the proof of 
Proposition~\ref{prop:1moment}, writing $W=W_0W_1$ in the notation of~\eqref{eq:WWW'}.
It follows from~\eqref{eq:W0} that 
$$
\prod_{p\leq M_{d,K}} p^{k(x)}\ll (\log x)^{C_1},
$$ 
where $C_1$  depends on the constant $C$ in the definition of $k(x)$, and on both $d$ and the number field $K$.
Similarly, 
$$
\prod_{p\mid h_\c} p^{k(x)}\leq h_\c^{k(x)\omega(h_\c)}\leq 
\exp\left(k(x)(\log h_\c)^2\right)\leq \exp\left(C_2(\log\log x)^3\right),
$$ 
since $h_\c\leq (\log x)^A$, 
for a constant $C_2$ that depends on $A$ and  
on the constant $C$ in the definition of $k(x)$. Hence it follows that 
\begin{equation}\label{eq:oat}
W_0\leq \exp\left(C'(\log\log x)^3\right),
\end{equation}
for a further constant $C'$ that depends on $d,K$ and $A,C$.


Since   $\gamma(W,f_\c(n))=\gamma(W_0,f_\c(n))\gamma(W_1,f_\c(n))$, we obtain
\begin{equation}\label{eq:S-to-T'}
\tilde N_{\c}(I)=
\sum_{\nu \bmod{W_0}} \gamma(W_0,f_\c(\nu))T(I;\nu), 
\end{equation}
where now
$$
T(I;\nu)=
\sum_{\substack{n\in I\\ 
n\equiv \nu\bmod{W_0}}}
\gamma(W_1,f_\c(n)).
$$
It follows from~\eqref{eq:calc} 
that 
$$
\gamma(W_1,f_\c(n))
=\prod_{p\mid P_\c(w)} \alpha_p
\prod_{\substack{p\mid P_\c(w)\\
p^{k(x)}\nmid f_\c(n)}}
r_K(p^{v_p(f_\c(n))}) 
\prod_{\substack{p\mid P_\c(w)\\
p^{k(x)}\mid f_\c(n)
}}
\alpha_p^{-1}\beta_{p^{k(x)}},
$$
where $\alpha_p$ and $ \beta_{p^{k(x)}}$ are given by~\eqref{eq:alpha-p} and ~\eqref{eq:beta-p}, 
respectively.
Hence
\begin{equation}\label{eq:pigeon}
T(I;\nu)=
U(I;\nu)\prod_{p\mid P_\c(w)} \alpha_p
\end{equation}
where
$$
U(I;\nu)=\sum_{\substack{n\in I
\\ n\equiv \nu\bmod{W_0}}} \tilde r_K(f_\c(n)_W),
$$
with the notation~\eqref{eq:r-tilde} for $\tilde r_K$ and where
$
n_W=\prod_{p^\nu \| n,~ p\mid P_\c(w)}p^\nu.
$
It remains to prove an asymptotic formula for $U(I;\nu)$.

 

 Note that $f_\c(n)\neq 0$ in $U(I,\nu)$, by hypothesis.
 Write 
 $$\mathcal{C}_w=\{n\in \mathbb{N}:\,\ p\mid n\implies p\mid P_\c(w)\}.$$
Let $d=f_\c(n)_W$.  Then there exist an integer $d'$ such that $f_\c(n)=dd'$, with 
$d\in \mathcal{C}_w$ and 
$\gcd(d',P_\c(w))=1$. 
Using M\"obius inversion to detect the coprimality condition, we obtain
\begin{align*}
U(I;\nu)
&=\sum_{\substack{k\in \mathcal{C}_w
}} g(k)
\sum_{\substack{n\in I
\\ n\equiv \nu\bmod{W_0}\\ 
k\mid f_\c(n)}} 1,
\end{align*}
where 
$
g=\tilde r_K* \mu$.
We will  break the outer  sum into two intervals, leading to 
\begin{equation}\label{eq:sum-sig123'}
U(I;\nu)=\Sigma_1+\Sigma_2,
\end{equation}
where
\begin{equation}\label{eq:Sigma1'}
\Sigma_1=
\sum_{\substack{k\in \mathcal{C}_w 
\\ k\leq x^{1/(2d)}}} g(k) 	
\sum_{\substack{n\in I
\\ n\equiv \nu\bmod{W_0}\\
k\mid f_\c(n)}} 1
\end{equation}
and 
$$
\Sigma_2=
\sum_{\substack{k\in \mathcal{C}_w 
\\ 
k>x^{1/(2d)}}} 
g(k) 	
\sum_{\substack{n\in I
\\ n\equiv \nu\bmod{W_0}\\
k\mid f_\c(n)}} 1.
$$


Beginning with the second sum, it follows from~\eqref{eq:size-g} and the inclusion  $I\subset (0,x]$, that 
\begin{align*}
\Sigma_2
&\ll
\sum_{\substack{k\in \mathcal{C}_w 
\\ 
k>x^{1/(2d)}}}
\tau(k)^{e+1}
\sum_{\substack{n\leq x\\ k\mid f_\c(n)}} 1.
\end{align*}
Suppose that $\omega(k)=r$. Then any $k$ in the sum satisfies
$$
x^{1/(2d)}<k=p_1^{j_1}\cdots p_{r}^{j_r}\leq w^{k(x)r},
$$
whence $\omega(f_\c(n))\geq \omega(k)=r > K_0$, with 
$$
K_0=\frac{\log x}{2 k(x)d\log w}=\frac{\sqrt{\log x}}{2k(x)d}.
$$
But then it follows that $1\leq 2^{-K_0} 2^{\omega(f_\c(n))}\leq 2^{-K_0} \tau(|f_\c(n)|)$, whence
$$
\Sigma_2
\ll
\frac{1}{2^{K_0}} 
\sum_{\substack{n\leq x}} \tau(|f_\c(n)|)^{e+2}.
$$
Since $f_\c$ has content  $h_\c\leq (\log x)^A$, 
it now follows from Lemma~\ref{lem:henriot} that 
\begin{equation}\label{eq:Sigma2-''}
\Sigma_2
\ll_A
\frac{x(\log x)^{K_1}}{2^{K_0}} 
\ll x\exp\left(-(\log x)^{1/4}\right),
\end{equation}
for a suitable constant $K_1$ only depending on $d,e$ and $A$.



Turning to the sum~\eqref{eq:Sigma1'}, 
it follows from~\eqref{eq:nature-g} that 
 $\gcd(k,W_0)=1$ 
and $v_p(k)\leq k(x)$, 
for any $k\in \mathcal{C}_w$ such that $g(k)\neq 0$. 
 Breaking into residue classes modulo $k$, we deduce that 
\begin{align*}
\Sigma_1
&=
\sum_{\substack{k\in \mathcal{C}_w\\ k\leq x^{1/(2d)}}} g(k) 	
\sum_{\substack{\alpha \bmod{k}\\ f_\c(\alpha)\equiv 0 \bmod{k}}}
\sum_{\substack{n\in I
\\ n\equiv \nu\bmod{W_0}\\
n\equiv \alpha\bmod{k}}}1\\
&=
\sum_{\substack{k\in \mathcal{C}_w\\ k\leq x^{1/(2d)}}} g(k) 	
\lambda_{f_\c}(k)
\left(\frac{|I|}{W_0k} +O(1)\right),
\end{align*}
where $\lambda_{f_\c}$ is given by~\eqref{eq:lambda}.
Since $f_\c$ has content  $h_\c\leq (\log x)^A$, 
 Lemma~\ref{lem:stewart} implies that $\lambda_{f_\c}(k)=O_\ve(k^{1-1/d+\ve}(\log x)^{A/d})$, for any $\ve>0$.
We have already seen in~\eqref{eq:size-g} that  $|g(k)|\leq \tau(k)^{e+1}=O_\ve(k^{\ve})$.
Hence
\begin{equation}\label{eq:Sigma1''}
\Sigma_1
=
\frac{|I|}{W_0}
\sum_{\substack{k\in \mathcal{C}_w\\ k\leq x^{1/(2d)}}} \frac{g(k) 	
\lambda_{f_\c}(k)}{k}
+O_{A,\ve}\left(x^{1-1/(2d)+\ve}\right).
\end{equation}
Since $|I|\leq x$, we see that 
$$
\frac{|I|}{W_0}
\sum_{\substack{k\in \mathcal{C}_w\\ k\leq x^{1/(2d)}}} \frac{g(k) 	
\lambda_{f_\c}(k)}{k}=
\frac{|I|}{W_0}
\sum_{\substack{k\in \mathcal{C}_w}} \frac{g(k) 	
\lambda_{f_\c}(k)}{k}+O\left( x\Upsilon(x)\right),
$$
where
\begin{equation}\label{eq:robin}
\Upsilon(x)=\sum_{\substack{k\in \mathcal{C}_w\\  k> x^{1/(2d)}}}  
 \frac{|g(k)| 
\lambda_{f_\c}(k)}{k}.
\end{equation}

We will show momentarily that 
\begin{equation}\label{eq:dove}
\Upsilon(x)\ll_A \exp\left(-\frac{\sqrt{\log x}}{16d}\right).
\end{equation}
From~\eqref{eq:dove} it will follow that 
$$
\Sigma_1
=\frac{|I|}{W_0}
\sum_{\substack{k\in \mathcal{C}_w}} \frac{g(k) 	
\lambda_{f_\c}(k)}{k}+O_A\left( x\exp\left(-\frac{\sqrt{\log x}}{16d}\right)
\right),
$$
in~\eqref{eq:Sigma1''}. Combining this with~\eqref{eq:Sigma2-''} in~\eqref{eq:sum-sig123'}, we obtain 
$$
U(I;\nu)
=\frac{|I|}{W_0}
\sum_{\substack{k\in \NN\\
p \mid  k \Rightarrow p\mid P_\c(w) }}
\hspace{-0.3cm}
 \frac{g(k) 	
\lambda_{f_\c}(k)}{k}+O_A\left( x\exp\left(-(\log x)^{1/4}\right)\right).
$$ 
We now insert this into~\eqref{eq:pigeon} and~\eqref{eq:S-to-T'}, noting that 
$c_\c(w)\leq \varphi^*(W)\ll \sqrt{\log x}$ and 
$$
\sum_{\nu \bmod{W_0}} \frac{\gamma(W_0,f_\c(\nu))}{W_0}=\sigma_{\mathbf{c}}(W_0),
$$
in the notation of ~\eqref{eq:sig-c}. 
The main term therefore agrees with the statement of the lemma.  In view of the upper bound 
\eqref{eq:oat} for $W_0$, we easily see that 
$$
c_\c(w)x\exp\left(-(\log x)^{1/4}\right)
\sum_{\nu \bmod{W_0}} \frac{\gamma(W_0,f_\c(\nu))}{W_0} 
\ll x\exp\left(-(\log x)^{1/8}\right).
$$
This 
thereby completes the proof of the lemma, subject to the verification of the  upper bound~\eqref{eq:dove} for $\Upsilon(x)$.

Returning to~\eqref{eq:robin}, we write $k=k_1k_2$ for coprime $k_1,k_2$, where $k_1$ is square-free and $k_2$ is square-full. 
Moreover, we may assume that $v_p(k_2)\leq k(x)$ for any prime $p\mid k_2$, since otherwise $g(k)=0$.
Furthermore any $p\mid k_1k_2$ is coprime to the content of $f_\c$, since $k\in \mathcal{C}_w$.
Note that  $\lambda_{f_\c}(k)=\lambda_{f_\c}(k_1)\lambda_{f_\c}(k_2)$ and 
$$
\lambda_{f_\c}(k_1)
\leq \prod_{p\| k_1}d
\leq d^{\omega(k_1)}\leq \tau(k_1)^d,
$$ 
by  Lemma~\ref{lem:stewart}.
Hence, on 
applying~\eqref{eq:size-g} to get 
 $|g(k)|\leq \tau(k)^{e+1}$, 
we deduce that 
$$
\Upsilon(x)\ll  (\log x)^A
\sum_{\substack{k_2\in \mathcal{C}_w\\
p^\nu\| k_2\Rightarrow 2\leq \nu\leq k(x)}}
 \frac{\tau(k_2)^{e+1}
\lambda_{f_\c}(k_2)}{k_2} 
\Sigma_{d+e+1}\left(\frac{x^{1/(2d)}}{k_2}\right),
$$
in the notation of~\eqref{eq:definition-sig0}.  Appealing to Lemma~\ref{lem:gouda} , we deduce that 
\begin{align*}
\Sigma_{d+e+1}\left(\frac{x^{1/(2d)}}{k_2}\right)
&\ll \exp\left(-\frac{\log x}{8d\log w}\right)
\exp\left(\frac{\log k_2}{4\log w}\right)\\
&
\leq \exp\left(-\frac{\sqrt{\log x}}{8d}\right) 
k_2^{1/(4\log w)},
\end{align*}
since 
$\log w=\sqrt{\log x}$.
Hence we obtain 
$$
\Upsilon(x)\ll 
 \exp\left(-\frac{\sqrt{\log x}}{8d}\right) (\log x)^{A} 
 \sum_{\substack{k_2\in \mathcal{C}_w\\
p^\nu\| k_2\Rightarrow 2\leq \nu\leq k(x)}}
 \frac{\tau(k_2)^{e+1}
\lambda_{f_\c}(k_2)}{k_2^{1-1/(4\log w)}
}.
$$
If $p\leq w$ then we have $p^{-1+1/(4\log w)}\leq \frac{3}{2}p^{-1}$. Hence it  follows from multiplicativity and Lemma~\ref{lem:stewart} that 
\begin{align*}
\sum_{\substack{k_2\in \mathcal{C}_w\\
p^\nu\| k_2\Rightarrow 2\leq \nu\leq k(x)}}
\hspace{-0.8cm}
& \frac{\tau(k_2)^{e+1}
\lambda_{f_\c}(k_2)}{k_2^{1-1/(4\log w)}}\\
&\leq \prod_{p\leq w}\left(1+\sum_{2\leq j\leq d} \frac{3d(j+1)^{e+1}}{2p} +\sum_{j>d} \frac{d(j+1)^{e+1}(3/2)^j}{p^{j/d}} \right)\\
&=\prod_{p\leq w}\left(1+O\left(\frac{1}{p}\right) \right)\\
&\ll (\log x)^B,
\end{align*}
for a constant $B$ depending only on $d$ and $e$. The desired upper bound~\eqref{eq:dove} for $\Upsilon(x) $ easily follows, which thereby completes the proof of the lemma. 
 \end{proof}

\end{document}
